% Document/LinearAlgebra/VectorSpaces/Foundations/OrderedBases.tex
% Section: Ordered Bases and Coordinates

\newpage
\section{Ordered Bases and Coordinates}

\begin{remark}[Motivation]
    A basis $B$ determines unique coefficients for each vector, but without an ordering
    on $B$, we cannot write these coefficients as a \emph{column vector}. Writing is inherently
    order-preserving, so to represent vectors in matrix form, we need \textbf{ordered bases}.
\end{remark}

\begin{definition}[Ordered Basis]\label{def:ordered-basis}
    An \textbf{ordered basis} of $V$ is a basis $B$ equipped with a \textbf{total order} $<$ on it.
    A \textbf{well-ordered basis} is the special case in which $<$ is a well-order. For finite
    $B$ the two notions coincide.
    When $B$ is finite, we write $\mathcal{B} = (b_0, \ldots, b_{n-1})$ as an $n$-tuple.
    When $B$ is infinite, we write $\mathcal{B} = (b_i)_{i \in I}$ for some totally ordered index set $I$.
\end{definition}

\begin{remark}[Total Order vs.\ Well-Order --- What Each Is For]\label{rem:order-buys}
    The order on a basis is \emph{auxiliary structure}, and the different layers of the theory
    need different amounts of it. Total order is the right default.
    \begin{itemize}
        \item \textbf{A total order suffices} for coordinates --- the unique finite expansion
              $v = \sum_{i \in F} a_i b_i$ already determines a well-defined coordinate vector
              (Definition~\ref{def:coordinates}) --- and for the \emph{entrywise} notion of
              triangularity, where one only compares a row index with a column index.
        \item \textbf{No order at all} is needed to define matrices or matrix multiplication:
              only \textbf{column-finiteness}. The matrix of a map and the product of two
              matrices are well-defined because each column has finite support, so each entry of
              a product is a support-finite sum over the contracted index --- never an ordered
              computation (see the chapter on Matrices).
        \item \textbf{Well-ordering} is needed only for transfinite \emph{constructions}:
              recovering an adapted basis from a complete flag (Section~\ref{sec:flags}),
              echelon and rank normal forms by recursion (Section~\ref{sec:rank-normal-form}),
              and Gram--Schmidt orthogonalisation (inner product spaces). These are recursions
              that require well-foundedness in order to terminate.
        \item \textbf{No generality is lost} by defaulting to a total order: by Zorn's lemma
              every basis admits a well-order, so requiring one excludes no space --- it merely
              fixes more structure than the lower layers use. Defaulting to a total order also
              keeps the order \emph{reversible}, which is what makes lower-triangularity statable
              (see the chapter on Endomorphisms).
    \end{itemize}
\end{remark}

\begin{remark}[Ordering Matters]
    The same underlying set with different orderings gives \emph{different} ordered bases.
    For example, $\{\CanonicalVector{0}, \CanonicalVector{1}\}$ in $\bb{R}^2$ yields two ordered bases: $(\CanonicalVector{0}, \CanonicalVector{1})$ and $(\CanonicalVector{1}, \CanonicalVector{0})$.
    These give different coordinate representations for the same vector.
\end{remark}

\begin{definition}[Coordinates]\label{def:coordinates}
    Let $\mathcal{B} = (b_i)_{i \in I}$ be an ordered basis for $V$ (so $I$ is totally ordered).
    For $v \in V$, the unique representation
    \[
        v = \sum_{i \in F} a_i b_i \quad (\text{$F \subseteq I$ finite, $a_i \neq 0$})
    \]
    determines the \textbf{coordinate vector} $\CoordVec[\mathcal{B}]{v} \in K^{(I)}$ defined by
    \[
        \CoordVec[\mathcal{B}]{v}(i) =
        \begin{cases}
            a_i & \text{if } i \in F \\
            0 & \text{otherwise}
        \end{cases}
    \]
    When $I$ is finite, we write $\CoordVec[\mathcal{B}]{v} = \begin{pmatrix} a_0 \\ \vdots \\ a_{n-1} \end{pmatrix}_{\mathcal{B}} \in K^n$.
\end{definition}

\begin{example}[Coordinate Dependence on Ordering]\label{ex:coord-ordering}
    In $\bb{R}^2$, let $v = (3, 2)$. With the standard basis:
    \begin{itemize}
        \item $\mathcal{B}_1 = (\CanonicalVector{0}, \CanonicalVector{1})$: $\CoordVec[\mathcal{B}_1]{v} = \begin{pmatrix} 3 \\ 2 \end{pmatrix}_{\mathcal{B}_1}$
        \item $\mathcal{B}_2 = (\CanonicalVector{1}, \CanonicalVector{0})$: $\CoordVec[\mathcal{B}_2]{v} = \begin{pmatrix} 2 \\ 3 \end{pmatrix}_{\mathcal{B}_2}$
    \end{itemize}
    The coordinates depend on the \textbf{ordering}, not just the set of basis vectors.
\end{example}

\begin{theorem}[Coordinate Bijection]\label{thm:coord-bijection}
    Let $\mathcal{B} = (b_i)_{i \in I}$ be an ordered basis for $V$. The coordinate map
    \[
        \Chart{\mathcal{B}}: V \to K^{(I)}, \quad v \mapsto \CoordVec[\mathcal{B}]{v}
    \]
    is a bijection. (It is in fact a linear isomorphism; see Theorem~\ref{thm:coord-iso}.)
\end{theorem}

\begin{proof}
    \textbf{Injectivity}: If $\CoordVec[\mathcal{B}]{v} = 0$, then all coefficients are zero, so $v = 0$.

    \textbf{Surjectivity}: For any $f \in K^{(I)}$ with finite support $F$, the vector
    $v = \sum_{i \in F} f(i) b_i$ satisfies $\CoordVec[\mathcal{B}]{v} = f$.
\end{proof}

\begin{remark}[Three Layers of a Vector]\label{rem:three-layers-vector}
    For a vector $v \in V$ with ordered basis $\mathcal{B} = (b_i)_{i \in I}$, we now have
    three distinct objects:
    \begin{enumerate}
        \item \textbf{The geometric object}: $v \in V$. Basis-free, lives in the abstract
              vector space.
        \item \textbf{The coordinate representation}: $\Chart{\mathcal{B}}(v) \in K^{(I)}$.
              An element of the canonical vector space $K^{(I)}$, obtained via the $K$-linear
              bijection $\Chart{\mathcal{B}}$. Basis-dependent, but still an algebraic object
              --- it is a vector, and $K$-linear maps can act on it.
        \item \textbf{The tuple of scalars}: the function $i \mapsto a_i$ where
              $v = \sum_{i \in I} a_i b_i$. This is the same underlying set-theoretic object
              as (2), but viewed purely as data --- a finitely-supported family of field
              elements, available for entrywise manipulation, printing, storage, and
              arithmetic.
    \end{enumerate}
    Layers (1) and (2) are related by the coordinate map $\Chart{\mathcal{B}}$, a $K$-linear
    bijection. Layers (2) and (3) are the same underlying set-theoretic object, but viewed
    with different intent: $K^{(I)}$ as a vector space (where addition and scalar
    multiplication are the relevant operations) versus $K^{(I)}$ as a set of functions
    $I \to K$ (where evaluation, support, and entrywise manipulation are the relevant
    operations). We identify them and rely on context to determine which structure is active,
    just as we identify $\bb{R}^2$ (vector space) with $\bb{R} \times \bb{R}$ (set of pairs)
    in everyday use.

    This pattern --- geometric object, coordinate representation in a canonical space,
    extracted data --- will recur at the level of linear maps in
    Section~\ref{sec:matrix-of-map}.
\end{remark}

\begin{proposition}[Canonical Basis of $K^{(I)}$]\label{prop:canonical-basis}
    For any set $I$, define $\CanonicalVector{i} \in K^{(I)}$ by $\CanonicalVector{i}(j) = \delta_{ij}$ (Kronecker delta).
    Then $(\CanonicalVector{i})_{i \in I}$ is a basis for $K^{(I)}$, called the \textbf{canonical basis}.
\end{proposition}

\begin{proof}
    \textbf{Generating}: For any $f \in K^{(I)}$ with finite support $F$:
    \[
        f = \sum_{i \in F} f(i) \cdot \CanonicalVector{i}
    \]
    (Verify: both sides agree at every $j \in I$.)

    \textbf{Linear independence}: Suppose $0 = \sum_{i \in F} a_i \CanonicalVector{i}$ for some finite $F$ with $a_i \neq 0$.
    Evaluating at any $j \in F$: $0 = a_j$, contradiction.
\end{proof}

\begin{remark}[Canonicity is Special to $K^{(I)}$]\label{rem:canonical-basis-special}
    The basis $(\CanonicalVector{i})_{i \in I}$ is \textbf{canonical} because it is defined intrinsically by
    the function-space structure of $K^{(I)}$ itself --- no arbitrary choice is involved.
    \textbf{An abstract $K$-vector space $V$ has no canonical basis}: any basis of $V$ requires a
    choice (equivalently, a choice of isomorphism $V \cong K^{(I)}$). Even $\bb{R}^n$ only has
    a ``standard'' basis because we keep track of the identification $\bb{R}^n = K^{(n)}$ where $n = \{0, 1, \ldots, n-1\}$.
    The canonical basis exists \emph{only} for spaces of the form $K^{(I)}$.
\end{remark}

\begin{corollary}[Dimension of $K^{(I)}$]\label{cor:dim-finite-support}
    $\Dim[K]{K^{(I)}} = \Card{I}$.
\end{corollary}

\begin{example}[Canonical Bases]\label{ex:canonical-bases}
    Spaces of the form $K^{(I)}$ have a canonical basis. The following are instances:
    \begin{itemize}
        \item $K^n = K^{(n)}$ where $n = \{0, 1, \ldots, n-1\}$: canonical basis $(\CanonicalVector{0}, \CanonicalVector{1}, \ldots, \CanonicalVector{n-1})$
        \item $K^{(\kappa)}$ for a cardinal $\kappa$: canonical basis $(e_\alpha)_{\alpha < \kappa}$
    \end{itemize}
    Other vector spaces have \emph{no} canonical basis. For instance, $K[X]$ has a natural ordered
    basis $(1, X, X^2, \ldots)$, but this is only ``natural'' because we implicitly choose the
    isomorphism $K[X] \cong K^{(\bb{N})}$ sending $X^n \mapsto \CanonicalVector{n}$. A different isomorphism would
    yield a different basis. The basis is not canonical --- only the isomorphism type is.
\end{example}

\begin{remark}[Cardinal Properties of $K^{(I)}$]
    We study the cardinal properties of $K^{(I)}$:
    \begin{enumerate}
        \item \textbf{Lower bound}: Embeddings $k \mapsto k \cdot \CanonicalVector{i_0}$ and $i \mapsto \CanonicalVector{i}$ give
              $\max(\Card{K}, \Card{I}) \leq \Card{K^{(I)}}$.
        \item \textbf{Bijection}: $K^{(I)} \sim \BigUnion[n \in \bb{N}] (K \setminus \{0\})^n \times \binom{I}{n}$ (disjoint union;
              a vector is determined by its finite support and nonzero values).
        \item \textbf{Finite case}: When both $K$ and $I$ are finite, $K^{(I)} = K^I$, so $\Card{K^{(I)}} = \Card{K}^{\Card{I}}$.
        \item \textbf{Infinite case}: When $\kappa = \max(\Card{K}, \Card{I}) \geq \aleph_0$:
              \[
                  \Card{K^{(I)}} = \sum_{n \in \bb{N}} \kappa^n = \aleph_0 \cdot \kappa = \kappa
              \]
    \end{enumerate}
\end{remark}
